\section{Conjugate Updating and Computational Mechanics}

When prior distributions belong to a family that, when multiplied by the likelihood, yields a posterior in the same functional family, the prior and likelihood are said to be \emph{conjugate}. Conjugacy allows exact analytical calculation without resorting to expensive Markov Chain Monte Carlo (MCMC) simulations.

\subsection{The Beta-Binomial Update Engine}

Consider a Bernoulli experiment observed \(n\) times with \(k\) successes. The likelihood is binomial:
\[
  p(k \mid \theta) = \binom{n}{k} \theta^k (1 - \theta)^{n - k}
\]
Placing a Beta prior on \(\theta \in [0, 1]\) with hyper-parameters \(\alpha_0, \beta_0 > 0\):
\[
  p(\theta) = \frac{\Gamma(\alpha_0 + \beta_0)}{\Gamma(\alpha_0)\Gamma(\beta_0)} \theta^{\alpha_0 - 1} (1 - \theta)^{\beta_0 - 1}
\]
Multiplying prior by likelihood yields the posterior kernel:
\[
  p(\theta \mid k) \propto \theta^{(\alpha_0 + k) - 1} (1 - \theta)^{(\beta_0 + n - k) - 1}
\]
Recognizing the functional form immediately identifies the posterior as:
\[
  \theta \mid k \sim \text{Beta}(\alpha_0 + k, \, \beta_0 + n - k)
\]

\begin{ExampleBox}{A/B Conversion Testing}
  Suppose a digital marketing system has a prior belief on click-through rate parameterized by \(\text{Beta}(2, 2)\). An in-market experiment runs across \(n = 100\) impressions and registers \(k = 18\) clicks.
  
  The updated parameters are:
  \[
    \alpha_{\text{post}} = 2 + 18 = 20, \quad \beta_{\text{post}} = 2 + 82 = 84
  \]
  The posterior expectation is:
  \[
    \mathbb{E}[\theta \mid k] = \frac{20}{20 + 84} = \frac{20}{104} \approx 0.1923
  \]
  Prior mean was \(0.50\), but the substantial evidence swiftly shifts mass toward the empirical rate of \(0.18\).
\end{ExampleBox}

\subsection{Algorithmic Implementation}

In production pipelines, Bayesian updates run as reproducible stateful functions:

\begin{HighlightedCodeBlock}[numbers=left]{Python}{4,5,6}
def update_beta_binomial(alpha_prior, beta_prior, n, k):
    """Computes exact conjugate posterior hyper-parameters."""
    alpha_post = alpha_prior + k
    beta_post = beta_prior + (n - k)
    mean_post = alpha_post / (alpha_post + beta_post)
    var_post = (alpha_post * beta_post) / (
        (alpha_post + beta_post)**2 * (alpha_post + beta_post + 1)
    )
    return alpha_post, beta_post, mean_post, var_post
\end{HighlightedCodeBlock}

\begin{figure}[htbp]
  \centering
  \begin{tikzpicture}[node distance=0.8cm, auto]
    \node[bayes/data] (prior) {\small Prior\\$\text{Beta}(\alpha_0, \beta_0)$};
    \node[bayes/data, right=1.2cm of prior] (data) {\small Observations\\$(n, k)$};
    \node[bayes/model, below=1.0cm of $(prior)!0.5!(data)$] (post) {\small Posterior\\$\text{Beta}(\alpha_n, \beta_n)$};
    \node[bayes/node, below=0.8cm of post] (decision) {\small Decision Policy};

    \draw[bayes/flow] (prior) -- (post);
    \draw[bayes/flow] (data) -- (post);
    \draw[bayes/flow] (post) -- (decision);
  \end{tikzpicture}
  \caption{Traceable computational flow from prior belief and evidence to decision boundary.}
\end{figure}

\begin{SummaryBox}{Variance Contraction}
  As the number of observed samples \(n \to \infty\), the variance of the posterior contracts at the rate \(\mathcal{O}(1/n)\). Epistemic uncertainty vanishes, and the posterior converges to a Dirac delta distribution centered on the true data-generating parameter.
\end{SummaryBox}
